\columnbreak
\section{Wechselrichter WR}
\subsection{Wirkungsgrade WR}

$
\boxed{
\eta_\text{MPPT} = \frac{P_\text{DC}}{P_\text{MPP}}
}$
$
\boxed{
\eta_\text{Um} = \frac{P_\text{AC}}{P_\text{DC}}
}$
$
\boxed{
\eta_\text{WR} = \eta_\text{MPPT} \cdot \eta_\text{Um}
}$
$
\boxed{
\eta_\text{WR} = \frac{P_\text{AC}}{P_\text{MPP}}
}
$

\vspace{0.1cm}

\begin{minipage}[t]{0.48\columnwidth}
    \renewcommand{\arraystretch}{1.2}
    \begin{tabular}{@{} l p{2.8cm} l @{}}
        $[\eta_\text{MPPT}]$ & Anpassungsw \dotfill & $\mathrm{-}$ \\
        $[\eta_\text{Um}]$   & Umwandlungsw \dotfill   & $\mathrm{-}$ \\
        $[\eta_\text{WR}]$   & Gesamtw \dotfill & $\mathrm{-}$ \\
    \end{tabular}
\end{minipage}
\hfill
\vrule width 1pt
\hfill
\begin{minipage}[t]{0.48\columnwidth}
    \renewcommand{\arraystretch}{1.2}
    \begin{tabular}{@{} l p{2.8cm} l @{}}
        $[P_\text{MPP}]$     & MPP-Leistung \dotfill     & $\mathrm{W}$ \\
        $[P_\text{DC}]$      & Leist. WR-Eingang \dotfill   & $\mathrm{W}$ \\
        $[P_\text{AC}]$      & Leist. WR-Ausgang \dotfill  & $\mathrm{W}$ \\
    \end{tabular}
\end{minipage}

\subsection{Dimmensionierung WR}

Typischer Wert: $SR_{\mathrm{AC}} = 1$, \quad Je nach Standort, AusRtg und Neigung: $ 0.8 \ldots 1.2$

\vspace{0.1cm}

$
\boxed{SR_{\mathrm{AC}} = \frac{P_{\mathrm{STC}}}{P_{\mathrm{AC\_N}}}}
$
$
\boxed{SR_{\mathrm{AC}} = \frac{n \cdot P_{\mathrm{STC,Modul}}}{P_{\mathrm{AC,N}}}}
$
$
\boxed{
n = \frac{SR_{\mathrm{AC}} \cdot P_{\mathrm{AC\_N}}}{P_{\mathrm{STC,Modul}}}
}
$

\vspace{0.1cm}

\renewcommand{\arraystretch}{1.2}
\begin{tabular}{@{} l p{6.5cm} l @{}}
    $[SR_{\mathrm{AC}}]$   & Auslegungsfaktor (Sizing Ratio) \dotfill                        & $-$ \\
    $[P_{\mathrm{STC}}]$   & Summe aller $P_{\mathrm{STC,Modul}}$ (bei STC) \dotfill         & $\mathrm{kW}$ \\
    $[P_{\mathrm{AC\_N}}]$  & Ausgangs-Nennleistung des Wechselrichter \dotfill                   & $\mathrm{kW}$ \\
    $[P_{\mathrm{STC,Modul}}]$   & Modulleistung bei STC \dotfill         & $\mathrm{kW}$ \\
    $[n]$   & Anzahl Module \dotfill                        & $-$ \\
\end{tabular}


\vspace{0.1cm}





\subsection{Spannungsdimensionierung von PV-Strings}

\subsubsection{Maximale Anzahl Module in einem String}

\begin{minipage}[t]{0.48\columnwidth}
\textbf{Aufgrund max zuläss. Spannung:}

$
\boxed{
n_{\mathrm{Max,L}} = \frac{U_{\mathrm{WR,Max}}}{U_{\mathrm{L,Modul \, (-10 \, ^\circ\mathrm{C})}}}
}
$
\end{minipage}
\hfill
\vrule width 1pt
\hfill
\begin{minipage}[t]{0.48\columnwidth}
\textbf{Aufgrund  MPP-Arbeitsbereichs:}

$
\boxed{
n_{\mathrm{Max,MPP}} = \frac{U_{\mathrm{WR,MPP,Max}}}{U_{\mathrm{MPP,Modul \, (-10 \, ^\circ\mathrm{C})}}}
}
$
\end{minipage}

\vspace{0.1cm}



\subsubsection{Minimale Anzahl Module in einem String}

\textbf{Aufgrund des MPP-Arbeitsbereichs:}

$
\boxed{
n_{\mathrm{Min,MPP}} = \frac{U_{\mathrm{WR,MPP,Min}}}{U_{\mathrm{MPP,Modul \, (70 \, ^\circ\mathrm{C})}}}
}
$

\vspace{0.1cm}

\renewcommand{\arraystretch}{1.2}
\begin{tabular}{@{} l p{6.5cm} l @{}}
    $[n_{\mathrm{Max,L}}]$        & Max. Anzahl Module (Spannungsgrenze) \dotfill & $-$ \\
    $[n_{\mathrm{Max,MPP}}]$      & Max. Anzahl Module (MPP-Spannungsbereich) \dotfill & $-$ \\
    $[n_{\mathrm{Min,MPP}}]$      & Min. Anzahl Module (MPP-Spannungsbereich) \dotfill & $-$ \\
    $[U_{\mathrm{WR,Max}}]$       & Max. Eingangsspannung des WR \dotfill & $\mathrm{V}$ \\
    $[U_{\mathrm{WR,MPP,Max}}]$   & Max. MPP-Spannungsbereich WR \dotfill & $\mathrm{V}$ \\
    $[U_{\mathrm{WR,MPP,Min}}]$   & Min. MPP-Spannungsbereich WR \dotfill & $\mathrm{V}$ \\
    $[U_{\mathrm{L,Modul \, (-10 \, ^\circ\mathrm{C})}}]$    & Leerlaufspannung Modul bei $\vartheta_{\mathrm{Zell}} =  -10 \, ^\circ \mathrm{C}$ \dotfill & $\mathrm{V}$ \\
    $[U_{\mathrm{MPP,Modul \, (70 \, ^\circ\mathrm{C})}}]$ & MPP-Spannung pro Modul bei Temperatur $\vartheta_{\mathrm{Zell}} = 70 \, ^\circ\mathrm{C}$ \dotfill & $\mathrm{V}$ \\
\end{tabular}



\subsection{Stromdimensionierung von PV-Strings}

\subsubsection{Maximale Anzahl paralleler Strings}

\begin{minipage}[t]{0.48\columnwidth}
\textbf{Aufgr. max zuläss. Kurzschlussstrom:}

$
\boxed{
n_{\mathrm{String,K}} = \frac{I_{\mathrm{WR,Max}}}{I_{\mathrm{K,Modul \, (70 \, ^\circ\mathrm{C})}}}
}
$
\end{minipage}
\hfill
\vrule width 1pt
\hfill
\begin{minipage}[t]{0.48\columnwidth}
\textbf{Aufgr. max Eingangsstrom:}

$
\boxed{
n_{\mathrm{String,MPP}} = \frac{I_{\mathrm{WR,Eingang,Max}}}{I_{\mathrm{MPP,Modul \, (70 \, ^\circ\mathrm{C})}}}
}
$
\end{minipage}

\vspace{0.1cm}

\renewcommand{\arraystretch}{1.2}
\begin{tabular}{@{} l p{6.5cm} l @{}}
    $[n_{\mathrm{String,K}}]$        & Max. Anzahl paralleler Strings (Kurzschlussgrenze) \dotfill & $-$ \\
    $[n_{\mathrm{String,MPP}}]$      & Max. Anzahl paralleler Strings (Eingangsstromgrenze) \dotfill & $-$ \\
    $[I_{\mathrm{WR,Max}}]$          & Max. PV-Kurzschlussstrom ($\mathrm{I_{SC\_PV}}$) \dotfill & $\mathrm{A}$ \\
    $[I_{\mathrm{WR,Eingang,Max}}]$  & Max. Eingangsstrom ($\mathrm{I_{DCmax}}$)\dotfill & $\mathrm{A}$ \\
    $[I_{\mathrm{K,Modul \, (70 \, ^\circ\mathrm{C})}}]$     & Kurzschlussstrom des Moduls bei $T = 70 \, ^\circ\mathrm{C}$ \dotfill & $\mathrm{A}$ \\
    $[I_{\mathrm{MPP,Modul \, (70 \, ^\circ\mathrm{C})}}]$   & MPP-Strom des Moduls bei $T = 70 \, ^\circ\mathrm{C}$ \dotfill & $\mathrm{A}$ \\
\end{tabular}


\subsection{PV-Netzkopplung}

\begin{tabular}{|l|p{3cm}|p{3.8cm}|}
\hline
\textbf{Aspekt} & \textbf{Früher: netzfolgend} & \textbf{Heute: netzführend}  \\
\hline
Frequenzhaltung & Abschaltung bei f > 50{,}2 Hz und f < 49{,}5 Hz & Kennlinie P(f) bei 50{,}2 bis 51{,}5 Hz  
 \\
\hline
Wirkleistung & Keine Regelung / Volleinspeisung & 
\begin{itemize}
  \item Wirkleistungsbegrenzung (z.\,B. 70\,\%)
  \item Wirkleistungsreduktion auf Anforderung durch Netzbetreiber (Rundsteuersignal)
\end{itemize} \\
\hline
Blindleistung & Cos~$\varphi$ = 1 & Cos~$\varphi$ -0{,}8 bis 0{,}8 (feste Einstellung, Rundsteuersignal, Q(U))  \\
\hline
Immunität & Abschaltung & Immunität gegen Netzstörung (Fault Ride Through)  \\
\hline
\end{tabular}

\vspace{0.1cm}

\textbf{Zukünftig: netzbildend}

\begin{minipage}[t]{0.48\columnwidth}
\begin{itemize}
  \item Systemspannung bilden
  \item Fehlerströme einspeisen
  \item Aktiv Oberwellen filtern
  \item Spannungsasymmetrien ausgleichen
\end{itemize}
\end{minipage}
\hfill
\begin{minipage}[t]{0.48\columnwidth}
\begin{itemize}
  \item Virtuelle Trägheit
  \item Erweitertes Fault Ride Through
  \item Netzstabilisierende Unterstützung
\end{itemize}
\end{minipage}